% DOC NO. RL-Z0-SCALING · REV. A
%
% This file IS the document. Edit it directly - ripdoc compiles it
% as it stands and stamps the provenance line at build time.
%
%   ripdoc build --only docs-scaling
\documentclass[fontpath=fonts/,logopath=logo/]{riposte-doc}
\usepackage{riposte-md}

\title{CHANNEL Z0 — Viewer Scaling (Station Bulletin № 6)}
\author{Riposte Laboratories}
\date{}
\ripdocno{RL-Z0-SCALING}
\riprev{A}
\riprunningtitle{CHANNEL Z0 — Viewer Scaling (Station Bulletin № 6)}

\begin{document}

\ripmakehead

\riptoc

\ripcodeopen
\ripcodesize{9.0}{13.95}
\begin{ripcodeblock}
┌───────────────────────────────────────────────────────────┐
│   CHANNEL Z0 · CH 0 · DESIG RL-Z0                         │
│   TRANSMITTER — VIEWER PATH                               │
│   hundreds, thousands, tens of thousands, millions        │
└───────────────────────────────────────────────────────────┘
\end{ripcodeblock}
\ripcodesizereset\ripcodeclose

What breaks first as the audience grows, and the one change that makes delivery cost stop depending on the viewer count. Playout (ErsatzTV) is unchanged at every tier; this bulletin is about the tower and the storefront.

\ripsection{\ripmdhead{The verdict}}

Move video delivery to object storage behind a CDN at the \textbf{hundreds} tier. After that, egress is free and the origin does one upload per segment regardless of how many people pull it. Every later tier is about the ingest path, the tower's side jobs (status, chat) and redundancy.

The reference build uses Cloudflare R2 (zero egress fee) fronted by Cloudflare's cache, because the storefront already lives on that account. Any S3-compatible store with a CDN in front works the same way.

\ripsection{\ripmdhead{The numbers that drive the plan}}

\begin{ripmdtable}{X[1.3,l]X[1.2,l]X[1.5,l]}
\ripmdth{FACT} & \ripmdth{VALUE} & \ripmdth{WHERE FROM} \\
Stream bitrate (passthrough) & \textasciitilde{}1.3 Mbps ≈ 0.6 GB/h per viewer & ErsatzTV 1200k + audio \\
One average concurrent viewer, per month & \textasciitilde{}430 GB & 0.6 GB/h × 720 h \\
Tower free egress & 10 TB/mo → \textasciitilde{}23 average viewers & Oracle Always Free \\
Oracle overage & \textasciitilde{}\$8.50/TB & Oracle price list, North America \\
Tower NIC & 480 Mbps → \textasciitilde{}350 viewers hard cap & E2.1.Micro shape \\
R2 egress & \$0 & Cloudflare R2 pricing \\
R2 writes (Class A) & \$4.50/M after 1M free; 3 s segments ≈ 1.7M/mo ≈ \$3 & Cloudflare R2 pricing \\
R2 reads (Class B) & \$0.36/M after 10M free; only on a CDN cache miss & Cloudflare R2 pricing \\
CDN edge TTL override floor, Free plan & 2 hours & Cloudflare cache docs \\
\end{ripmdtable}

\ripsubsection{\ripmdhead{If nothing changes}}

\begin{ripmdtable}{X[1,l]X[1,l]X[1,l]X[1,l]}
\ripmdth{AVG CONCURRENT} & \ripmdth{EGRESS / MO} & \ripmdth{ORACLE BILL} & \ripmdth{FITS THE NIC?} \\
23 & 10 TB & \$0 & yes \\
100 & 43 TB & \textasciitilde{}\$280 & yes, peaks under 350 \\
1,000 & 430 TB & \textasciitilde{}\$3,600 & no \\
10,000+ & 4 PB+ & no & no \\
\end{ripmdtable}

\ripsubsection{\ripmdhead{With R2 behind the CDN}}

\begin{ripmdtable}{X[2,l]X[1,l]}
\ripmdth{ITEM} & \ripmdth{COST / MO} \\
Segment + playlist writes & \textasciitilde{}\$3 \\
Storage (rolling few minutes) & \textasciitilde{}\$0 \\
Origin reads with cached playlists and Smart Tiered Cache & under \$5 \\
Origin reads if playlists are \riphi{not} cached & \$0.31 per average viewer (\$310 at 1,000) \\
\end{ripmdtable}

\textbf{Playlist caching is the single lever.} Owncast uploads every \ripmono{.m3u8} with \ripmono{Cache-Control: no-cache, no-store, must-revalidate} (hard-coded in its S3 storage code), and a CDN never caches \ripmono{no-store}. On Cloudflare's Free plan the edge-TTL override floor is 2 hours, useless for a live playlist. The fix is a \textbf{Cache Response Rule} (available on every plan, 10 rules on Free): strip \ripmono{no-store} and \ripmono{no-cache}, set \ripmono{s-maxage=2} on \ripmono{*.m3u8} for the HLS host. Segments are fine as-is: Owncast sends \ripmono{max-age=N} on \ripmono{.ts} and a Cache Rule marking the host eligible for cache does the rest.

If a Cache Response Rule ever stops being an option, the fallback is a sidecar on the tower that watches Owncast's local HLS directory and uploads with its own headers. Owncast's S3 code is then unused. That works with any plan and any CDN.

\ripsection{\ripmdhead{Tier 0: now (under about 25 average viewers)}}

No delivery change. Two chores that every later tier needs:

\begin{ripbullets}
\item Alerting on off-air. The sentinel's \ripmono{z0-tower-live} check runs hourly; make sure a FAIL reaches a human the same hour, not in the next morning's brief.
\item Queue management (SQM/cake) on the home WAN so torrent traffic cannot starve the uplink. It has taken the stream down twice.
\end{ripbullets}

\ripsection{\ripmdhead{Tier 1: hundreds}}

Cost: about \$5–10/mo. Effort: one weekend.

\begin{ripbullets}
\item Create the bucket (\ripmono{ch0-hls}), attach the custom domain \ripmono{hls.ch0.ripostelabs.xyz}, add a lifecycle rule "delete after 1 day" as a backstop. Owncast deletes its own playback window; the rule catches crashes.
\item Cache Rules on the HLS host: eligible for cache, Smart Tiered Cache on. Cache Response Rule for playlists as above. Confirm \ripmono{cf-cache-status: HIT} on both \ripmono{.ts} and \ripmono{.m3u8} from two regions before switching the player.
\item Owncast: Storage → S3 with the R2 endpoint, bucket, key pair, region \ripmono{auto}, path-style on, and the video serving endpoint set to the custom domain. A freshly enabled R2 account answers the S3 endpoint with a TLS handshake failure (alert 40) for some minutes until its certificate exists; every upload fails until then, so probe the endpoint with \ripmono{curl} before switching. \ripmono{tools/z0-r2-offload.sh} applies all of it through the admin API. Then \riphi{restart the uplink} on the playout host: Owncast only builds the new pipeline when a new RTMP session starts, and a config change alone changes nothing.
\item Storefront: nothing to change. Owncast keeps serving the master playlist from the watch host and rewrites it so the variant playlist and every segment come from the HLS host; the player follows. One tiny request per viewer session still touches the tower, which is fine until tier 2. The R2 domain must answer with \ripmono{Access-Control-Allow-Origin: *}; the bucket's CORS rule provides it.
\item Raise Owncast's latency level to 3. Expect 20–40 s glass-to-glass. The NOW SHOWING title will lead the picture by about that much.
\item Viewer count: Owncast no longer sees segment requests, so its own count reads 0. Take unique visitors on the HLS host from the CDN's analytics API over a 5-minute window and feed the storefront chip from that.
\end{ripbullets}

What stays: the tower, RTMP ingest, the single rendition.

\ripsection{\ripmdhead{Tier 2: thousands}}

The video path is unchanged. The tower's side jobs become the limit.

\begin{ripbullets}
\item \ripmono{/api/status} is polled every 30 s per viewer: 1,000 viewers is 33 req/s, 10,000 is 333 req/s on a one-core box. Split hostnames: an ingest host (DNS-only, RTMP) and the watch host (proxied, JSON only). Caddy already stamps \ripmono{Cache-Control: public, s-maxage=10} on \ripmono{/api/status}; the CDN honours origin \ripmono{max-age} on every plan, so the tower sees one request per 10 s per cache tier. Block \ripmono{/hls/*} on the proxied host so no video crosses the CDN outside the object store. Cloudflare's terms require a paid product (R2 counts) for video through the CDN.
\item Chat: Owncast's chat is one websocket per viewer. Turn it off. The storefront's comment relay is the chat.
\item Storefront Worker: static assets are free and unlimited, but the comment relay's Worker requests count against 100k/day on the free plan. Move to Workers Paid before a busy night.
\item Second rendition (480p for phones) needs CPU at the packager: the free ARM shape if the retry timer ever wins, or a paid 2-core shape (\textasciitilde{}\$35/mo). A single 1.3 Mbps rendition is acceptable until people complain.
\item Standby tower: a second Owncast with the same stream key. \ripmono{z0-leader.sh} already stands the uplink down on an unhealthy target; give it a second target so it fails over instead of stopping.
\end{ripbullets}

\ripsection{\ripmdhead{Tier 3: tens of thousands}}

Delivery cost is still flat. Design changes:

\begin{ripbullets}
\item Take Owncast out of the video path. Package at home: ffmpeg reads the ErsatzTV channel, makes a 3-rung ladder with the playout host's hardware encoder (1080p/720p/360p) and writes HLS locally; a small uploader puts segments and playlists in the bucket with the right headers. No cloud transcoder bill. The tower keeps only status and comments, or retires.
\item Ingest redundancy: a second packager in the cloud fed by the same RTMP push, writing to a second bucket prefix. The master playlist lists both, so players fail over on their own.
\item A paid CDN plan for analytics, more rules and support.
\item Watch cache hit ratio and Class B reads daily. A falling hit ratio is the early warning that cost is about to scale with viewers again.
\end{ripbullets}

\ripsection{\ripmdhead{Tier 4: millions}}

The architecture does not change: object store plus CDN does not cost more per viewer. What changes is the relationship with the CDN.

\begin{ripbullets}
\item Around 1 Tbps sustained from a self-serve zone triggers a conversation. Budget for an enterprise plan, or be ready with a second CDN pulling from the same bucket. The origin is plain HTTP objects, so multi-CDN is a DNS change, not a rebuild.
\item Never route this through a per-minute video service. At \$1 per 1,000 delivered minutes, one million viewers cost \textasciitilde{}\$60k per hour.
\item Two independent ingest sites (home and cloud), each with its own packager and bucket prefix. The home WAN is otherwise the single point of failure.
\item Comments: one coordinator object for the whole channel serialises every write. Shard by minute or room before this tier.
\item Viewer counting moves fully to CDN analytics. Nothing per-viewer touches an origin.
\end{ripbullets}

\end{document}
